\begin{resumo}[RESUMO]
\vspace*{-6mm}

Operar uma planta de processo envolve mais do que sintonizar malhas: é preciso escolher o ponto de operação, decidir o que manter constante diante dos distúrbios, controlar e, por fim, acompanhar se a planta cumpre o que foi decidido. Este trabalho trata dessa última etapa, a supervisão, e chama de política de operação o conjunto de decisões que a planta deve cumprir. O objetivo é mostrar que uma camada supervisória construída sobre o Kubernetes consegue acompanhar uma planta a partir de uma política declarada pelo usuário, sem que o supervisor conheça a planta, usando o Processo Tennessee Eastman (TEP) como planta e o OPC~UA como meio de comunicação. Como o trabalho se insere na ênfase de Integração de Sistemas de Controle, o foco está em ligar as camadas, e não em sintonizar controladores. O modelo de Downs e Vogel foi reescrito em Rust como uma planta que roda de forma contínua e publica os seus sinais por OPC~UA. Um \textit{historian} em Python coleta esses sinais e calcula estatísticas por janela de tempo e um índice de desempenho das malhas de controle. Um controlador próprio, escrito em Go e executado num \textit{cluster} Kubernetes local, lê três recursos declarados pelo usuário (a função de custo, a política de operação e a planta), avalia a planta contra a política do Modo~1 a cada 30~s e grava o veredito no próprio Kubernetes, de onde uma IHM web o lê. A planta reproduziu o caso base de Downs e Vogel, com custo de 170,63~\$/h contra 170,61~\$/h publicados, e, na versão com a água de resfriamento integrada, o ótimo do Modo~1 de Ricker, com 114,32~\$/h contra 114,31~\$/h. Com a política fixa, o veredito ficou em conformidade na operação normal e acusou a degradação causada pelo distúrbio IDV(6) 1,15~h depois dele, na ordem física dos acontecimentos: primeiro a vazão de produto, depois o custo. O índice das malhas, sozinho, não percebeu o distúrbio, e foi preciso acrescentar um critério de seguimento do \textit{setpoint}. Uma releitura dos dados mostrou que o índice estava calculado com a escala invertida; na escala do artigo original, ele teria acusado o distúrbio antes do veredito econômico, mas também acusa transitórios inofensivos. Num terceiro experimento, com malhas PI e sem malha de nível do reator, a planta se perdeu em cerca de 33~h enquanto o custo caía, e só o limite de nível declarado na política percebeu, e tarde. Conclui-se que o objetivo foi atingido, com limites: o mecanismo de estado desejado e estado observado do Kubernetes se mostrou adequado para a supervisão, mas uma política útil precisa olhar também para os inventários e para as tendências, e não só para o custo e para as malhas. Fazer o supervisor agir sobre a planta fica para trabalhos futuros.

\textbf{Palavras-chave}: Integração de sistemas de controle. Processo Tennessee Eastman. Supervisão de processos. Política de operação. Kubernetes. OPC UA.
\end{resumo}
